\documentclass[12pt]{article}
\usepackage[a4paper,margin=1in]{geometry}
\usepackage[T1]{fontenc}
\usepackage{lmodern,microtype}
\usepackage{amsmath,amssymb,amsthm,mathtools}
\usepackage{booktabs,array,enumitem,xcolor}
\usepackage[numbers,sort&compress]{natbib}
\usepackage[colorlinks=true,linkcolor=blue!55!black,citecolor=blue!55!black,urlcolor=blue!55!black]{hyperref}
\linespread{1.04}

\newtheorem{theorem}{Theorem}[section]
\newtheorem{proposition}[theorem]{Proposition}
\newtheorem{corollary}[theorem]{Corollary}
\theoremstyle{definition}
\newtheorem{definition}[theorem]{Definition}
\theoremstyle{remark}
\newtheorem{remark}[theorem]{Remark}
\newcommand{\cX}{\mathcal X}
\newcommand{\cY}{\mathcal Y}
\newcommand{\cK}{\mathcal K}
\newcommand{\cH}{\mathcal H}
\newcommand{\Ran}{\operatorname{Ran}}
\newcommand{\rank}{\operatorname{rank}}
\newcommand{\norm}[1]{\left\lVert#1\right\rVert}

\title{Canonical Polar Realization of Reset-Memory Packets\\
\large An Exact Finite-History Solution to the Source-Space Type Problem}
\author{Liang Wang}
\date{July 2026}

\begin{document}
\maketitle

\begin{abstract}
The reset atlas of RH-151 selects a clock-rank packet from a recursive
source-memory Gram matrix.  The packet therefore lives in source-column
space, whereas the moving-cloud determinant of RH-80 lives in a transfer
space.  RH-162 correctly identified the resulting ambient-space type gap,
but did not exploit the special Gram structure of the archived reset data.

We factor that structure exactly.  If $S_k=A^kS_0$ and
$\widehat S_k=S_k/\norm{S_k}_F$, define
\[
 F_t x=(\widehat S_t x,\sqrt\eta\widehat S_{t-1}x,
 \ldots,\eta^{t/2}\widehat S_0x).
\]
Then the recursive memory is precisely $M_t=F_t^*F_t$.  For any rank-$r$
spectral packet $P_t$ on which $M_t$ is positive, the partial polar map
\[
 J_t=F_tP_t\bigl(P_tM_tP_t|_{\Ran P_t}\bigr)^{-1/2}
\]
is an isometry from $\Ran P_t$ into the normalized history space.  Its range
projection is canonical, independent of packet-frame gauge and covariant
under every unitary change of source coordinates.  In an eigenframe
$U_t^*M_tU_t=\Lambda_t$, the formula becomes
$V_t=F_tU_t\Lambda_t^{-1/2}$.

This closes a finite-dimensional source-memory-to-history realization,
denoted $X_{\rm mem\to hist}$.  It does not identify the history space with
the RH-80 transfer/determinant space and therefore does not close the
physical interface $X_{\rm phys}$.  A 192-case complex synthetic audit
checks the Gram, polar, source-gauge, and packet-gauge identities with maximum
residual below $5.5\times10^{-15}$.  The theorem is exact; the audit validates
only its implementation.  No all-level limit, Riesz cloud, canonical
determinant, Hilbert--Polya operator, or Riemann-hypothesis conclusion is
claimed.
\end{abstract}

\noindent\textbf{Keywords:} Gram factorization; polar decomposition; reset
packet; finite history; gauge covariance; nonselfadjoint dynamics.

\section{The precise type problem}

RH-151 constructs, at each frozen scale and time, the top clock-rank spectral
packet of a positive memory matrix \citep{WangRH151}.  The memory acts on the
columns of the source $S_0$, not on the state space on which the production
transfer operator acts.  Writing both packets with the same letter suppresses
a genuine type distinction:
\[
 P_t^{\rm reset}\in\mathcal B(\cX),
 \qquad
 \Pi_t^{\rm cloud}\in\mathcal B(\cH_{\rm transfer}).
\]
RH-162 therefore required an ambient realization before applying a
packet-to-Riesz theorem \citep{WangRH162}.

The first question is narrower than the full physical bridge.  Does the
particular reset memory already carry a canonical Hilbert-space realization,
without fitting target eigenvectors or choosing an arbitrary injection?  The
answer is yes.  The target is a normalized finite-history space determined by
the source dynamics itself.  This paper constructs that target and proves its
canonicity.  Whether it is the correct transfer target is left open.

\section{Normalized source memory}

Let $\cX$ and $\cY$ be finite-dimensional complex Hilbert spaces.  Let
$A\in\mathcal B(\cY)$, let $S_0:\cX\to\cY$ be nonzero, and suppose
$S_k=A^kS_0\ne0$ for $0\le k\le t$.  Fix $0\le\eta<1$ and put
\[
 \alpha_k=\norm{S_k}_F,
 \qquad \widehat S_k=\alpha_k^{-1}S_k.
\]
The RH-94/RH-151 memory recursion is
\begin{equation}\label{eq:memory-recursion}
 M_{-1}=0,
 \qquad
 M_k=\widehat S_k^*\widehat S_k+\eta M_{k-1}.
\end{equation}
Every $M_k$ is positive semidefinite and
$\operatorname{tr}M_k=1+\eta+\cdots+\eta^k$.

Define the history space
\[
 \cK_t=\underbrace{\cY\oplus\cdots\oplus\cY}_{t+1\text{ copies}}
\]
with current time first, and define $F_t:\cX\to\cK_t$ by
\begin{equation}\label{eq:history-factor}
 F_tx=
 \bigl(\widehat S_tx,\sqrt\eta\widehat S_{t-1}x,
 \ldots,\eta^{t/2}\widehat S_0x\bigr).
\end{equation}

\begin{theorem}[Exact history Gram factorization]\label{thm:gram}
For every finite time $t$,
\begin{equation}\label{eq:gram-factorization}
 M_t=F_t^*F_t
 =\sum_{k=0}^t\eta^{t-k}\widehat S_k^*\widehat S_k.
\end{equation}
Consequently $\ker M_t=\ker F_t$ and
$\rank M_t=\rank F_t$.
\end{theorem}

\begin{proof}
The direct-sum norm has no cross terms, hence
\[
 \langle F_tx,F_ty\rangle
 =\sum_{k=0}^t\eta^{t-k}
 \langle\widehat S_kx,\widehat S_ky\rangle.
\]
This proves the second identity in \eqref{eq:gram-factorization}.  Iterating
\eqref{eq:memory-recursion} gives the same sum.  Finally,
$\langle M_tx,x\rangle=\norm{F_tx}^2$, so the kernels agree; rank equality
follows in finite dimension.
\end{proof}

The factor $F_t$ is not introduced by an eigenvector fit.  It is forced by
the actual recursion, including the normalization and forgetting weight.
Thus the source memory already has an intrinsic realization, albeit on a
history space rather than the desired transfer space.

\section{Canonical realization of a reset packet}

Let $P$ be a rank-$r$ orthogonal projection on $\cX$.  Assume
\begin{equation}\label{eq:positive-packet}
 PM_tP|_{\Ran P}>0.
\end{equation}
This condition is automatic when $P$ is spanned by $r$ positive eigenvalues
of $M_t$.  On $\Ran P$ define
\begin{equation}\label{eq:partial-polar}
 J_{t,P}=F_tP\bigl(PM_tP|_{\Ran P}\bigr)^{-1/2}.
\end{equation}

\begin{theorem}[Canonical packet isometry]\label{thm:polar}
Under \eqref{eq:positive-packet}, $J_{t,P}:\Ran P\to\cK_t$ is an isometry:
\[
 J_{t,P}^*J_{t,P}=I_{\Ran P}.
\]
Its range is $F_t(\Ran P)$, and its ambient projection is
\begin{equation}\label{eq:realized-projection}
 Q_{t,P}=J_{t,P}J_{t,P}^*.
\end{equation}
Both $J_{t,P}$ and $Q_{t,P}$ are determined by $(F_t,P)$; no target
eigenbasis is chosen.
\end{theorem}

\begin{proof}
By Theorem~\ref{thm:gram},
\[
 (F_tP)^*(F_tP)|_{\Ran P}=PM_tP|_{\Ran P}.
\]
Substitution in \eqref{eq:partial-polar} gives
$J_{t,P}^*J_{t,P}=I$.  The inverse square root is invertible on $\Ran P$,
so multiplying by it does not change the image of $F_tP$.  Uniqueness is the
uniqueness of the polar partial isometry on the initial space
$(\ker F_tP)^\perp=\Ran P$ \citep{Higham2008}.
\end{proof}

Suppose now that $U:\mathbb C^r\to\cX$ is an orthonormal eigenframe,
\[
 U^*U=I_r,
 \qquad P=UU^*,
 \qquad U^*M_tU=\Lambda=\operatorname{diag}(\lambda_1,\ldots,\lambda_r),
\]
where every $\lambda_j>0$.  Then Theorem~\ref{thm:polar} becomes
\begin{equation}\label{eq:frame-formula}
 V=F_tU\Lambda^{-1/2},
 \qquad V^*V=I_r,
 \qquad Q_{t,P}=VV^*.
\end{equation}
This is the formula used by the numerical audit.

\section{Gauge covariance and degeneracy}

There are two unrelated gauges: a basis change in source space and a basis
change inside the selected packet.  A canonical construction must handle
both.

\begin{proposition}[Source-coordinate covariance]\label{prop:source-gauge}
Let $R$ be unitary on $\cX$ and replace every source state by
$S_k'=S_kR$.  Then
\[
 F_t'=F_tR,
 \qquad M_t'=R^*M_tR.
\]
If $P'=R^*PR$, the realized history projection is unchanged:
\[
 Q'_{t,P'}=Q_{t,P}.
\]
\end{proposition}

\begin{proof}
Frobenius norms are invariant under right unitary multiplication, so every
normalized block transforms by $\widehat S_k'=\widehat S_kR$.  This proves
the first two identities.  The positive compression transforms by unitary
conjugacy, and functional calculus gives
$(P'M_t'P')^{-1/2}=R^*(PM_tP)^{-1/2}R$ on the transformed range.  Hence
$J'_{t,P'}R^*=J_{t,P}$ and the range projections agree.
\end{proof}

\begin{proposition}[Packet-frame equivariance]\label{prop:packet-gauge}
In \eqref{eq:frame-formula}, replace $U$ by $UW$ for a unitary
$W\in\mathbb C^{r\times r}$.  If the positive factor is transformed rather
than rediagonalized, the polar frame becomes $VW$.  Therefore $VV^*$ is
independent of packet frame.
\end{proposition}

At a repeated eigenvalue, an eigenframe is not canonical, but the spectral
projection is canonical whenever the selected cluster is separated from its
complement.  Theorem~\ref{thm:polar} is deliberately formulated with $P$,
not with an ordered eigenbasis.  It therefore survives internal degeneracy.
If the rank boundary itself has zero gap, the choice of $P$ is not canonical;
this is a packet-selection issue already separated in RH-151.

\section{Conditioning and stable evaluation}

Formula \eqref{eq:frame-formula} is exact but can be a poor floating-point
algorithm when $\lambda_r$ is very small.  Direct division by
$\sqrt{\lambda_r}$ magnifies the Gram and eigenvector errors.  A stable
implementation forms $X=F_tU$ and computes its thin singular value
decomposition
\[
 X=L\Sigma R^*.
\]
The polar isometry is then
\begin{equation}\label{eq:svd-polar}
 V=LR^*,
 \qquad H=R\Sigma R^*,
 \qquad X=VH.
\end{equation}
In exact arithmetic \eqref{eq:svd-polar} equals
\eqref{eq:frame-formula}.  Numerically it enforces orthonormal columns without
explicitly applying a large inverse square root \citep{Higham2008}.

\begin{remark}[What ``canonical'' means numerically]
An SVD routine may choose different singular-vector phases or rotate a
repeated singular subspace.  Those choices alter the frame $V$ but not the
projection $VV^*$.  The mathematically canonical object is the polar partial
isometry, or its range projection when only the realized subspace is needed.
\end{remark}

\section{Synthetic implementation audit}

The reproducibility script generates 192 complex cases with source dimensions
$6,9,12$, history lengths $1,2,4,7$, ranks one and three, and eight random
trials per configuration.  Each case applies an independent unitary source
gauge and an independent packet-frame gauge.  The maximum residuals are:
\begin{center}
\begin{tabular}{lr}
\toprule
identity & maximum residual\\
\midrule
$F_t^*F_t=M_t$ & $7.16\times10^{-16}$\\
$V^*V=I$ & $2.02\times10^{-15}$\\
polar factorization $F_tU=VH$ & $1.42\times10^{-15}$\\
source-gauge covariance & $1.26\times10^{-15}$\\
packet-gauge equivariance & $5.43\times10^{-15}$\\
packet-gauge projection invariance & $2.20\times10^{-15}$\\
\bottomrule
\end{tabular}
\end{center}
The audit is ordinary floating-point arithmetic and is not an interval
certificate.  It checks that the code implements the proved identities.  It
does not validate the upstream physical matrices or establish an asymptotic
bound.

\section{A universal property for competing ambient realizations}

The history factor is canonical because the memory recursion specifies it,
but the Gram matrix by itself never determines an ambient Hilbert space
absolutely.  The precise residual freedom is unitary equivalence.

\begin{proposition}[Universal property of the packet Gram]
\label{prop:universal}
Let $G:\Ran P\to\mathcal Z$ be another injective realization on a Hilbert
space $\mathcal Z$ satisfying
\begin{equation}\label{eq:competing-gram}
 G^*G=PM_tP|_{\Ran P}.
\end{equation}
Let
\[
 W_G=G(G^*G)^{-1/2}
\]
be its polar isometry.  Then there is a unique unitary
\[
 U_G:F_t(\Ran P)\longrightarrow G(\Ran P)
\]
such that
\begin{equation}\label{eq:universal-unitary}
 U_GJ_{t,P}=W_G.
\end{equation}
Conversely, every unitary $U_G$ on the realized packet range produces a
realization with the same compressed Gram.
\end{proposition}

\begin{proof}
Both $J_{t,P}$ and $W_G$ are isometries with initial space $\Ran P$.
Define $U_G(J_{t,P}x)=W_Gx$.  Isometry of both maps makes this definition
well posed and inner-product preserving, and surjectivity onto $G(\Ran P)$
is immediate.  Uniqueness follows because $J_{t,P}(\Ran P)$ is the whole
domain of $U_G$.  Conversely, if
$G=U_GJ_{t,P}(PM_tP)^{1/2}$, then $G^*G=PM_tP$.
\end{proof}

This proposition gives a necessary exact test for a proposed transfer-space
realization.  If a candidate analysis map $G_t$ on the transfer side has the
same packet Gram as $F_t$, then the packet portions of history and transfer
space are canonically unitarily equivalent through
\eqref{eq:universal-unitary}.  The remaining problem is no longer an
equal-rank identification; it is to construct $G_t$ from the physical
transfer dynamics and verify \eqref{eq:competing-gram}, or a controlled
approximate version of it.

\begin{corollary}[No exact Gram match, no exact isometric bridge]
If a proposed transfer realization $G$ satisfies
$G^*G\ne PM_tP$ on $\Ran P$, then no unitary between the two realized packet
ranges can intertwine both analysis maps while fixing their source
coordinates.  Any bridge must include a nontrivial positive polar correction.
\end{corollary}

Thus RH-172 supplies not only one ambient realization but also a comparison
principle for every future one.  The physical leaf
$X_{\rm hist\to transfer}$ can be attacked by estimating the Gram mismatch,
the induced polar correction, and the commutator with time evolution.

\section{Exact progress and remaining interface}

The physical leaf $X_{\rm phys}$ from RH-171 can now be refined as
\begin{equation}\label{eq:x-split}
 X_{\rm phys}
 =X_{\rm mem\to hist}\wedge X_{\rm hist\to transfer}.
\end{equation}
This paper proves the first factor at every finite time for every nonzero
selected eigenvalue.  The second factor remains open.  In particular, the
history space $\cK_t$ grows with $t$, is not the RH-80 transfer space, and has
not been shown to intertwine the noisy transfer operator, its square, or the
moving determinant cloud.

The construction nevertheless removes one ambiguity.  Future attempts may
no longer identify a source packet with a transfer packet by notation; they
must either:
\begin{enumerate}[leftmargin=2em]
\item construct a bounded map from the canonical history realization into
the transfer space and control its polar correction and commutator; or
\item bypass reset memory and construct the Riesz cloud from a different
intrinsic model.
\end{enumerate}
Failure of the first option would reject only the reset-history branch.

\section{Theorem boundary}

The exact results are finite-history Gram factorization, canonical polar
realization, source-coordinate covariance, packet-gauge invariance, and the
stable SVD evaluation rule.  They establish $X_{\rm mem\to hist}$.

They do not establish a history-to-transfer intertwiner, a physical contour
or Riesz projection, a uniform all-level packet gap, a Schatten complement,
a canonical relative determinant, directed marked-trace convergence, a
self-adjoint completion, a $T\log T$ law, a prime-power trace formula, a zeta
divisor identity, or the Riemann Hypothesis.

\bibliographystyle{plainnat}
\bibliography{references}
\end{document}
